% Generated by Claude Opus 5 (Anthropic) — compressed §2 and §3 for the 4-page limit.
% §2: 1278 -> ~400 words.  §3: 999 -> ~545 words.

\section{Methodology}
\label{sec:method}

FISH builds a model of a deployed ROS~2 application as the ROS~2 and CUDA
runtimes execute it. Four kinds of observation, on one clock, suffice for this.
\emph{Object identity}: every \texttt{rcl} handle with its creation and
destruction, so that callbacks are attributed to the right node and executor
even when handles are reused. \emph{Dispatch}: the start and end of every
callback execution, whatever its trigger. \emph{Data flow}: which callback
publishes to which topic and which callback consumes it, including deliveries
that never leave the process. \emph{GPU work}: every runtime API call and the
kernels, copies, memsets and synchronisations it produces.

\noindent
\textbf{Instrumentation:}
FISH builds on the handle-matched tracepoints of
\texttt{ros2\_tracing}~\cite{bedard-2022-ros2tracing}, which come in two phases
--- \emph{initialisation} events, recording every object as it is created, and
\emph{execution} events, recording what then happens to it --- and enables 42
events: 14 stock and 28 added. The additions fall into six groups. Six finalize
events close the object lifetimes that the stock \texttt{*\_init} events open,
and a seventh binds a timer to its node and period atomically; together they let
a callback be attributed to the right object even after a handle is reused.
Callback start and end at \texttt{execute\_waitable} and
\texttt{execute\_client}, with two client round-trip and four action-server
events, complete the five executable kinds of the \texttt{rclcpp} executor.
Three link events, emitted at the publish, receive and client-request sites,
turn a publish into an edge from the callback that issued it. The intra-process
subscription waitable and two NITROS links cover the deliveries that never reach
\texttt{rcl}. Five registration events give Python nodes the same callback chain
as \texttt{rclcpp}, except for waitables. Four introspection events record the
executor's type and thread count and each callback's group membership. GPU work
is recorded by CUPTI through \texttt{nsys} with a minimal profile and aligned to
the LTTng clock at trace ingest. Correlation ids then bind each GPU record to
its API call, and each API call to the callback window that issued it.

\noindent
\textbf{Acquisition:}
No application source modification is required: the instrumented ROS~2 libraries
are installed on the host or in a containerised environment, and the
\texttt{ros2} entry point is replaced by the FISH entry point, which wraps the
original command. Because CUPTI observes only processes started under the
profiler, the entry point first walks the launch description in describe-only
mode and launches the containers that will host GPU components directly under
\texttt{nsys}. This must happen at spawn time: launch loads composable nodes
into a container only once, so a container restarted later would come up empty.
Standalone GPU processes are handled by the FISH daemon, which polls
\texttt{/proc} and restarts a process under the profiler when it detects a CUDA
context. A system-stable signal then gates the session: the daemon declares the
system up when the node list has been unchanged over five polls and all
components are loaded, and only then starts the input: a rosbag replay, or for a
live system a stimulus schedule, a list of timed commands issued at fixed
offsets from the system-stable instant. Phase boundaries are derived from the
trace itself, so that comparisons use the steady state regardless of how the
input arrived: the initialisation phase ends when every periodic timer has fired
once, the warm-up phase when every recurring callback has completed its first
execution.

\noindent
\textbf{From trace to model:}
Data ingest places both traces in one time base. The model generator builds the
hierarchy process--executor--node--entity--callback, attaches the interaction
edges of Section~\ref{sec:model}, and for every callback that submits GPU work
derives a typed DAG per invocation.





\section{Model Definition}
\label{sec:model}

The FISH model follows the layers of abstraction described in
Section~\ref{sec:intro}, plus two that ROS~2 does not name: a host, which holds
the whole deployment, and a layer for the node interfaces, which we call
entities. The model is a hierarchical graph $G=(V,E)$ with an expansion operator
$\mathcal{X}$.

\noindent
\textbf{Vertices:}
Each vertex is a tuple $v=(t_v,\mathit{id}_v,A_v,Z_v,\mathit{ns}_v)$ --- type,
identifier, attributes, children in the next layer, and namespace --- whose
per-layer specification is given in Fig.~\ref{fig:model}a. The children of one
layer partition the layer below it. An entity is one of the node interfaces of
Section~\ref{sec:intro}, the thing the executor waits on; a callback is what it
then runs. A publisher is waited on by nobody, and a callback of any kind may
publish, so it is not a vertex but an aspect of the entity whose callback
publishes through it. Namespaces are not a layer; $\mathit{ns}_v$ only groups
vertices for presentation.

\noindent
\textbf{Edges:}
An edge is an interaction, or a set of interactions, between two vertices: a
tuple $(\mathit{src},\mathit{dst},t,m)$ of its endpoints, the interaction type
and its measurement. There are five types (Table~\ref{tab:edges}), each
recovered from one of the two sources of Section~\ref{sec:method}: the
initialisation events, which declare how the application is wired, and the
execution events, which show what it did.

\begin{table}[t]
\centering
\caption{The five kinds of interaction an edge can stand for.}
\label{tab:edges}
\resizebox{\columnwidth}{!}{%
\begin{tabular}{lll}
\hline
$t$ & What it models & $m$ \\ \hline
$\mathit{msg}$            & \begin{tabular}[c]{@{}l@{}}a message from a publisher\\ to a subscription\end{tabular}      & hop latency \\ \hline
$\mathit{dep}$            & \begin{tabular}[c]{@{}l@{}}a service request and\\ its response\end{tabular}                 & round trip  \\ \hline
$\mathit{state\text{-}polled}$   & \begin{tabular}[c]{@{}l@{}}a message pulled when\\ needed rather than dispatched\end{tabular} & data age    \\ \hline
$\mathit{state\text{-}inferred}$ & \begin{tabular}[c]{@{}l@{}}a value one callback stores\\ and another later reads\end{tabular} & data age    \\ \hline
$\mathit{join}$           & inputs that complete a set                                                                   & ---         \\ \hline
\end{tabular}%
}
\end{table}

\noindent
$\mathit{msg}$ and $\mathit{dep}$ come from initialisation, and are laid down
first at the entity layer: a $\mathit{msg}$ edge wherever a publisher and a
subscription share a topic, a $\mathit{dep}$ pair wherever a client and a server
share a service. At this point every edge in the graph runs between entities.
Execution then does two things to each. It finds the callback that actually
issued the interaction --- registration already names the subscription's
callback, but any callback of the publishing node may publish, so only the trace
says which one did --- and it copies the edge down to the callback layer, where
it takes its measurements. The remaining three come from execution alone, and
are constructed directly at the callback layer. Figure~\ref{fig:model}c gives
the trace pattern each is decided from.

\noindent
Every edge is then lifted along the containment. Writing $Z^{*}_{v}$ for the
callbacks under $v$, an edge between any two vertices is the union of the
callback-layer edges between their callbacks,
%
\[
  E(u,v)=\bigcup_{a\in Z^{*}_{u},\ b\in Z^{*}_{v}} E(a,b),
\]
%
so two callbacks exchanging a message are also two entities exchanging it, two
nodes, two executors. The edges marked \textcircled{\scriptsize 2} in
Fig.~\ref{fig:model}b are one and the same message: carried from the entity that
declared it down to the callback that issued it, and drawn again at every layer
above, the same interaction seen from further away.

\noindent
$\mathcal{X}$ is the inverse step: it replaces a vertex by the vertices of its
$Z$ set, and its edges by theirs. Applying it to one vertex and not another is
how the graph is read at mixed granularity --- one executor opened into its
nodes while the rest stay whole.

\noindent
\textbf{Correctness:}
A map $\chi$ takes every vertex-creating ROS~2 API call to the trace records it
must emit and the vertices it must create --- a node constructor yields
$1N+7E+7F$, because it also creates the six parameter services and the
time-source subscription --- and FISH scans the sources of every binary its
traced processes mapped to sum the expected records and vertices. Correctness is
then the two containments $\chi_r(A)\subseteq R$ and $\chi_v(A)\subseteq V$ over
the executed calls $A$: instrumentation completeness, every expected record is
in the trace, and extraction completeness, every expected vertex is in the
graph.
